\begin{register}{H}{SDHOST\_CTRL\_REG}{0x0000}\label{regdesc:SDHOSTCTRLREG}
\regfield{(reserved)}{6}{26}{000000}%
\regfield{SDHOST\_USE\_INTERNAL\_DMAC}{1}{25}{{0}}%
\regfield{(reserved)}{13}{12}{0000000000000}%
\regfield{SDHOST\_CEATA\_DEVICE\_INTERRUPT\_STATUS}{1}{11}{{0}}%
\regfield{SDHOST\_SEND\_AUTO\_STOP\_CCSD}{1}{10}{{0}}%
\regfield{SDHOST\_SEND\_CCSD}{1}{9}{{0}}%
\regfield{SDHOST\_ABORT\_READ\_DATA}{1}{8}{{0}}%
\regfield{SDHOST\_SEND\_IRQ\_RESPONSE}{1}{7}{{0}}%
\regfield{SDHOST\_READ\_WAIT}{1}{6}{{0}}%
\regfield{(reserved)}{1}{5}{0}%
\regfield{SDHOST\_INT\_ENABLE}{1}{4}{{0}}%
\regfield{(reserved)}{1}{3}{0}%
\regfield{SDHOST\_DMA\_RESET}{1}{2}{{0}}%
\regfield{SDHOST\_FIFO\_RESET}{1}{1}{{0}}%
\regfield{SDHOST\_CONTROLLER\_RESET}{1}{0}{{0}}%
\reglabel{Reset}\regnewline%
\begin{regdesc}\begin{reglist}
\label{fielddesc:SDHOSTCONTROLLERRESET}\item [SDHOST\_CONTROLLER\_RESET] 配置是否复位控制器。此位在控制器复位后自动清除。\\0：无效\\1：复位\\(R/W)
\label{fielddesc:SDHOSTFIFORESET}\item [SDHOST\_FIFO\_RESET] 配置是否复位 FIFO。此位在 FIFO 复位后自动清除。\\0：无效\\1：复位\\(R/W)
\label{fielddesc:SDHOSTDMARESET}\item [SDHOST\_DMA\_RESET] 配置是否复位 DMA 接口。此位在 DMA 复位后自动清除。\\0：无效\\1：复位\\(R/W)
\label{fielddesc:SDHOSTINTENABLE}\item [SDHOST\_INT\_ENABLE] 写入 1 使能全局中断端口。 (R/W)
\label{fielddesc:SDHOSTREADWAIT}\item [SDHOST\_READ\_WAIT] 配置是否向 SDIO 卡发送读取等待。\\0：清除读取等待\\1：生成读取等待\\(R/W)

\item [见下页...]
\end{reglist}\end{regdesc}
\end{register}

\addtocounter{Regfloat}{-1}
\begin{register}{H}{SDHOST\_CTRL\_REG}{0x0000}
\begin{regdesc}\begin{reglist}
\item [接上页...]

\label{fielddesc:SDHOSTSENDIRQRESPONSE}\item [SDHOST\_SEND\_IRQ\_RESPONSE] 配置是否发送自动中断请求 (IRQ) 响应。此位在响应发送后自动清除。\\
主机等待 MMC 卡中断时，会发出 CMD40 命令，并等待来自 MMC 卡的中断响应。如果主机希望 SD/MMC 退出等待中断状态，可以将此位置 1，此时 SD/MMC 命令状态机在总线上发送 CMD40 响应并返回到空闲状态。\\
0：无效\\1：发送自动 IRQ 响应\\(R/W)
\label{fielddesc:SDHOSTABORTREADDATA}\item [SDHOST\_ABORT\_READ\_DATA] 配置是否中止读取数据。此位在数据状态机重置为空闲后自动清除。\\
在读取操作期间发出暂停命令后，软件轮询卡获取暂停事件发生的时间。暂停事件发生后，软件将设置该位，该位将重置正在等待下一个数据块的数据状态机。\\
0：无效\\1：中止读取数据\\(R/W)
\label{fielddesc:SDHOSTSENDCCSD}\item [SDHOST\_SEND\_CCSD] 配置是否向 CE-ATA 设备发送命令完成信号禁用 (CCSD)。CCSD 发送给设备后，SD/MMC 会自动清除 SDHOST\_SEND\_CCSD 位，并将 SDHOST\_RINTSTS\_REG 寄存器中的命令完成 (CD) 位置 1。\\
只有在当前命令期待 CCS（即，RW\_BLK），并且为 CE-ATA 设备使能了中断时，软件才会设置此位。\\
0:无效\\1：向 CE-ATA 设备发送 CCSD\\(R/W)
\label{fielddesc:SDHOSTSENDAUTOSTOPCCSD}\item [SDHOST\_SEND\_AUTO\_STOP\_CCSD] 配置是否向 CE-ATA 设备发送内部生成的 STOP 命令 (CMD12)。发送 CCSD 后，SD/MMC 会自动清除 SDHOST\_SEND\_AUTO\_STOP\_CCSD 位。发送此内部生成的 STOP 命令后，SDHOST\_RINTSTS\_REG 中的自动命令完成 (ACD) 位会被置 1。\\
SDHOST\_SEND\_AUTO\_STOP\_CCSD 和 SDHOST\_SEND\_CCSD 位需一起配置。 \\
0:无效\\1：向 CE-ATA 设备发送内部生成的 STOP 命令 (CMD12)\\(R/W)
\label{fielddesc:SDHOSTCEATADEVICEINTERRUPTSTATUS}\item [SDHOST\_CEATA\_DEVICE\_INTERRUPT\_STATUS] 写入 1 以在 CE-ATA 设备中使能中断（ATE 控制寄存器中的 nIEN = 0）。 \\
在电源复位或 CE-ATA 设备的任何其他复位后，软件应适当地写入此位。复位后，CE-ATA 设备的中断通常被禁用（nIEN = 1）。如果主机使能了 CE-ATA 设备的中断，那么软件应设置此位。 (R/W)
\label{fielddesc:SDHOSTUSEINTERNALDMAC}\item [SDHOST\_USE\_INTERNAL\_DMAC] 配置是否使用内部 DMA 进行数据传输。\\0：不使用\\1：使用\\(R/W)
\end{reglist}\end{regdesc}
\end{register}


\begin{register}{H}{SDHOST\_CLKDIV\_REG}{0x0008}\label{regdesc:SDHOSTCLKDIVREG}
\regfield{SDHOST\_CLK\_DIVIDER3}{8}{24}{{0x0}}%
\regfield{SDHOST\_CLK\_DIVIDER2}{8}{16}{{0x0}}%
\regfield{SDHOST\_CLK\_DIVIDER1}{8}{8}{{0x0}}%
\regfield{SDHOST\_CLK\_DIVIDER0}{8}{0}{{0x0}}%
\reglabel{Reset}\regnewline%
\begin{regdesc}\begin{reglist}
\label{fielddesc:SDHOSTCLKDIVIDER0}\item [SDHOST\_CLK\_DIVIDER0] 配置时钟分频器 0 分频系数的二分之一。该字段为 0 时分频系数为 1。 (R/W)
\label{fielddesc:SDHOSTCLKDIVIDER1}\item [SDHOST\_CLK\_DIVIDER1] 配置时钟分频器 1 分频系数的二分之一。该字段为 0 时分频系数为 1。 (R/W)
\label{fielddesc:SDHOSTCLKDIVIDER2}\item [SDHOST\_CLK\_DIVIDER2] 配置时钟分频器 2 分频系数的二分之一。该字段为 0 时分频系数为 1。 (R/W)
\label{fielddesc:SDHOSTCLKDIVIDER3}\item [SDHOST\_CLK\_DIVIDER3] 配置时钟分频器 3 分频系数的二分之一。该字段为 0 时分频系数为 1。 (R/W)
\end{reglist}\end{regdesc}
\end{register}


\begin{register}{H}{SDHOST\_CLKSRC\_REG}{0x000C}\label{regdesc:SDHOSTCLKSRCREG}
\regfield{(reserved)}{28}{4}{0000000000000000000000000000}%
\regfield{SDHOST\_CLKSRC\_REG}{4}{0}{{0x0}}%
\reglabel{Reset}\regnewline%
\begin{regdesc}\begin{reglist}
\label{fielddesc:SDHOSTCLKSRCREG}\item [SDHOST\_CLKSRC\_REG] 配置两个 SD 卡输出到 sdhost\_cclk\_out[1:0] 信号的时钟分频器源。每张卡两位，位 [1:0] 对应卡 0，位 [3:2] 对应卡 1。：\\
0：时钟分频器 0\\
1：时钟分频器 1\\
2：时钟分频器 2\\
3：时钟分频器 3\\ (R/W)
\end{reglist}\end{regdesc}
\end{register}


\begin{register}{H}{SDHOST\_CLKENA\_REG}{0x0010}\label{regdesc:SDHOSTCLKENAREG}
\regfield{(reserved)}{14}{18}{00000000000000}%
\regfield{SDHOST\_LP\_ENABLE}{2}{16}{{0x0}}%
\regfield{(reserved)}{14}{2}{00000000000000}%
\regfield{SDHOST\_CCLK\_ENABLE}{2}{0}{{0x0}}%
\reglabel{Reset}\regnewline%
\begin{regdesc}\begin{reglist}
\label{fielddesc:SDHOSTCCLKENABLE}\item [SDHOST\_CCLK\_ENABLE] 配置是否使能两个 SD 卡时钟和一个 MMC 卡时钟。每张卡一位，位 0 对应卡 0。\\
0：时钟禁用\\
1：时钟使能\\ (R/W)
\label{fielddesc:SDHOSTLPENABLE}\item [SDHOST\_LP\_ENABLE] 配置卡处于空闲状态时是否禁用时钟。每张卡一位，位 0 对应卡 0。\\
0：时钟禁用\\
1：时钟使能\\ (R/W)
\end{reglist}\end{regdesc}
\end{register}


\begin{register}{H}{SDHOST\_TMOUT\_REG}{0x0014}\label{regdesc:SDHOSTTMOUTREG}
\regfield{SDHOST\_DATA\_TIMEOUT}{24}{8}{{0xffffff}}%
\regfield{SDHOST\_RESPONSE\_TIMEOUT}{8}{0}{{0x40}}%
\reglabel{Reset}\regnewline%
\begin{regdesc}\begin{reglist}
\label{fielddesc:SDHOSTRESPONSETIMEOUT}\item [SDHOST\_RESPONSE\_TIMEOUT] 配置响应超时时间。单位：卡输出时钟周期。 (R/W)
\label{fielddesc:SDHOSTDATATIMEOUT}\item [SDHOST\_DATA\_TIMEOUT] 配置卡数据读取超时时间和主机填充数据超时时间。超时计数器在卡时钟停止后才开始计数。单位：卡输出时钟周期。\\
注意：如果超时时间约为 100 毫秒，需使用软件定时器。此时需要禁用读取数据超时中断 (DRTO)。\\ (R/W)
\end{reglist}\end{regdesc}
\end{register}


\begin{register}{H}{SDHOST\_CTYPE\_REG}{0x0018}\label{regdesc:SDHOSTCTYPEREG}
\regfield{(reserved)}{14}{18}{00000000000000}%
\regfield{SDHOST\_CARD\_WIDTH8}{2}{16}{{0x0}}%
\regfield{(reserved)}{14}{2}{00000000000000}%
\regfield{SDHOST\_CARD\_WIDTH4}{2}{0}{{0x0}}%
\reglabel{Reset}\regnewline%
\begin{regdesc}\begin{reglist}
\label{fielddesc:SDHOSTCARDWIDTH4}\item [SDHOST\_CARD\_WIDTH4] 配置卡是否为 1 位模式或 4 位模式。每张卡一位，位 0 对应卡 0。\\
0：1 位模式\\
1：4 位模式\\
 (R/W)
\label{fielddesc:SDHOSTCARDWIDTH8}\item [SDHOST\_CARD\_WIDTH8] 配置卡是否为 8 位模式。每张卡一位，位 0 对应卡 0。\\
0：非 8 位模式\\
1：8 位模式\\
 (R/W)
\end{reglist}\end{regdesc}
\end{register}


\begin{register}{H}{SDHOST\_BLKSIZ\_REG}{0x001C}\label{regdesc:SDHOSTBLKSIZREG}
\regfield{(reserved)}{16}{16}{0000000000000000}%
\regfield{SDHOST\_BLOCK\_SIZE}{16}{0}{{0x200}}%
\reglabel{Reset}\regnewline%
\begin{regdesc}\begin{reglist}
\label{fielddesc:SDHOSTBLOCKSIZE}\item [SDHOST\_BLOCK\_SIZE] 配置块大小。 (R/W)
\end{reglist}\end{regdesc}
\end{register}


\begin{register}{H}{SDHOST\_BYTCNT\_REG}{0x0020}\label{regdesc:SDHOSTBYTCNTREG}
\regfield{SDHOST\_BYTE\_COUNT}{32}{0}{{0x000200}}%
\reglabel{Reset}\regnewline%
\begin{regdesc}\begin{reglist}
\label{fielddesc:SDHOSTBYTECOUNT}\item [SDHOST\_BYTE\_COUNT] 配置待传输的字节数，应为块传输的块大小的整数倍。\\
对于未定义字节长度的数据传输，该字段应设置为 0。此时，主机需发送停止/中止命令终止数据传输。 (R/W)
\end{reglist}\end{regdesc}
\end{register}


\begin{register}{H}{SDHOST\_CMDARG\_REG}{0x0028}\label{regdesc:SDHOSTCMDARGREG}
\regfield{SDHOST\_CMDARG\_REG}{32}{0}{{0x000000}}%
\reglabel{Reset}\regnewline%
\begin{regdesc}\begin{reglist}
\label{fielddesc:SDHOSTCMDARGREG}\item [SDHOST\_CMDARG\_REG] 配置要传递给卡的命令参数。 (R/W)
\end{reglist}\end{regdesc}
\end{register}


\begin{register}{H}{SDHOST\_CMD\_REG}{0x002C}\label{regdesc:SDHOSTCMDREG}
\regfield{SDHOST\_START\_CMD}{1}{31}{{0}}%
\regfield{(reserved)}{1}{30}{0}%
\regfield{SDHOST\_USE\_HOLE\_REG}{1}{29}{{1}}%
\regfield{SDHOST\_VOLT\_SWITCH}{1}{28}{{0}}%
\regfield{(reserved)}{4}{24}{0000}%
\regfield{SDHOST\_CCS\_EXPECTED}{1}{23}{{0}}%
\regfield{SDHOST\_READ\_CEATA\_DEVICE}{1}{22}{{0}}%
\regfield{SDHOST\_UPDATE\_CLOCK\_REGISTERS\_ONLY}{1}{21}{{0}}%
\regfield{SDHOST\_CARD\_NUMBER}{5}{16}{{0x0}}%
\regfield{SDHOST\_SEND\_INITIALIZATION}{1}{15}{{0}}%
\regfield{SDHOST\_STOP\_ABORT\_CMD}{1}{14}{{0}}%
\regfield{SDHOST\_WAIT\_PRVDATA\_COMPLETE}{1}{13}{{0}}%
\regfield{SDHOST\_SEND\_AUTO\_STOP}{1}{12}{{0}}%
\regfield{SDHOST\_TRANSFER\_MODE}{1}{11}{{0}}%
\regfield{SDHOST\_READ\_WRITE}{1}{10}{{0}}%
\regfield{SDHOST\_DATA\_EXPECTED}{1}{9}{{0}}%
\regfield{SDHOST\_CHECK\_RESPONSE\_CRC}{1}{8}{{0}}%
\regfield{SDHOST\_RESPONSE\_LENGTH}{1}{7}{{0}}%
\regfield{SDHOST\_RESPONSE\_EXPECT}{1}{6}{{0}}%
\regfield{SDHOST\_CMD\_INDEX}{6}{0}{{0x0}}%
\reglabel{Reset}\regnewline%
\begin{regdesc}\begin{reglist}
\label{fielddesc:SDHOSTCMDINDEX}\item [SDHOST\_CMD\_INDEX] 配置发送到卡的命令索引编号。 (R/W)
\label{fielddesc:SDHOSTRESPONSEEXPECT}\item [SDHOST\_RESPONSE\_EXPECT] 配置是否需要卡响应。\\0：不需要卡响应\\1：需要卡响应\\(R/W)
\label{fielddesc:SDHOSTRESPONSELENGTH}\item [SDHOST\_RESPONSE\_LENGTH] 配置卡的响应长度。\\0：短响应\\1：长响应\\(R/W)
\label{fielddesc:SDHOSTCHECKRESPONSECRC}\item [SDHOST\_CHECK\_RESPONSE\_CRC] 配置是否检查响应 CRC。\\0：不检查\\1：检查响应 CRC\\(R/W)
\label{fielddesc:SDHOSTDATAEXPECTED}\item [SDHOST\_DATA\_EXPECTED] 配置是否需要传输数据。\\0：不需要数据传输\\1：需要数据传输\\(R/W)
\label{fielddesc:SDHOSTREADWRITE}\item [SDHOST\_READ\_WRITE] 配置数据传输方向。\\0：从卡读取\\1：写入卡\\(R/W)
\label{fielddesc:SDHOSTTRANSFERMODE}\item [SDHOST\_TRANSFER\_MODE] 配置数据传输模式。\\0：传输数据块\\1：传输数据流\\(R/W)
\item [见下页...]
\end{reglist}\end{regdesc}
\end{register}

\addtocounter{Regfloat}{-1}
\begin{register}{H}{SDHOST\_CMD\_REG}{0x002C}
\begin{regdesc}\begin{reglist}
\item [接上页...]
\label{fielddesc:SDHOSTSENDAUTOSTOP}\item [SDHOST\_SEND\_AUTO\_STOP] 配置数据传输结束时是否自动发送停止命令。\\0：不发送\\1：发送\\(R/W)
\label{fielddesc:SDHOSTWAITPRVDATACOMPLETE}\item [SDHOST\_WAIT\_PRVDATA\_COMPLETE] 配置是否立即发送命令。\\0：即使前一个数据传输尚未完成，也立即发送命令\\1：在发送命令之前，等待前一个数据传输完成\\ (R/W)
\label{fielddesc:SDHOSTSTOPABORTCMD}\item [SDHOST\_STOP\_ABORT\_CMD] 配置停止/中止命令是否可以停止数据传输。\\0：停止或中止命令都不能停止当前的数据传输。如果当前选择的卡设备，之前已向它发送过中止命令，或它不在数据传输模式中，则此位应设置为 0\\
1：停止或中止命令旨在停止当前正在进行的数据传输\\ (R/W)
\label{fielddesc:SDHOSTSENDINITIALIZATION}\item [SDHOST\_SEND\_INITIALIZATION] 配置是否在发送此命令之前发送初始化序列（80 个时钟周期的 1），以便控制器在向卡发送命令之前初始化时钟。上电后，必须先向卡发送初始化序列，然后才能发送命令。所以向卡发送第一个命令时需要置位该字段。 \\
0：不发送初始化序列\\
1：发送初始化序列\\
 (R/W)
\label{fielddesc:SDHOSTCARDNUMBER}\item [SDHOST\_CARD\_NUMBER] 配置正在使用的卡号和正在访问的卡的物理插槽号。在 SD-only 模式下，支持最多两张卡。  (R/W)
\label{fielddesc:SDHOSTUPDATECLOCKREGISTERSONLY}\item [SDHOST\_UPDATE\_CLOCK\_REGISTERS\_ONLY] 配置是否仅将 SDHOST\_CLKDIV\_REG、SDHOST\_CLRSRC\_REG 和 SDHOST\_CLKENA\_REG 时钟寄存器的值更新到卡时钟域。 \\
0：正常发送命令序列\\
1：不发送命令，只更新时钟寄存器值到卡时钟域\\
该位可以确保更改时钟频率或停止时钟时无需向卡发送命令，也不会触发命令完成中断。 (R/W)
\label{fielddesc:SDHOSTREADCEATADEVICE}\item [SDHOST\_READ\_CEATA\_DEVICE] 配置主机是否在读取 CE-ATA 设备，用于开启或禁用读取超时指示。\\
0：主机未读取（RW\_REG 或 RW\_BLK）CE-ATA 设备，开启数据读取超时指示\\
1：主机正在读取（RW\_REG 或 RW\_BLK）向 CE-ATA 设备，禁用数据读取超时指示\\
 (R/W)
\item [见下页...]
\end{reglist}\end{regdesc}
\end{register}

\addtocounter{Regfloat}{-1}
\begin{register}{H}{SDHOST\_CMD\_REG}{0x002C}
\begin{regdesc}\begin{reglist}
\item [接上页...]
\label{fielddesc:SDHOSTCCSEXPECTED}\item [SDHOST\_CCS\_EXPECTED] 配置是否需要命令完成信号 (CCS)。\\
0：CE-ATA 设备中未使能中断（ATA 控制寄存器中的 nIEN = 1），或命令不需要 CE-ATA 设备发送 CCS\\
1：CE-ATA 设备中使能了中断（nIEN = 0），并且 RW\_BLK 命令需要 CE-ATA 设备发送 CCS\\
 (R/W)
\label{fielddesc:SDHOSTVOLTSWITCH}\item [SDHOST\_VOLT\_SWITCH] 配置是否使能电压切换。\\
0：没有电压切换\\
1：使能电压切换。\\
CMD11 中此位必须为 1。(R/W)
\label{fielddesc:SDHOSTUSEHOLEREG}\item [SDHOST\_USE\_HOLE\_REG] 配置是否使用保持寄存器 (hold register)。\\
0：CMD 和 DATA 绕过保持寄存器发送到卡\\
1：CMD 和 DATA 通过保持寄存器发送到卡\\  (R/W)
\label{fielddesc:SDHOSTSTARTCMD}\item [SDHOST\_START\_CMD] 配置是否开始命令。\\
0：不开始\\
1：开始\\
一旦命令由 CIU 服务，此位将自动清除。\\ 此位为 1 时，主机不能尝试写入任何命令寄存器，否则会触发硬件锁定错误。 (R/W/SC)
\end{reglist}\end{regdesc}
\end{register}


\begin{register}{H}{SDHOST\_FIFOTH\_REG}{0x004C}\label{regdesc:SDHOSTFIFOTHREG}
\regfield{(reserved)}{1}{31}{0}%
\regfield{SDHOST\_DMA\_MULTIPLE\_TRANSACTION\_SIZE}{3}{28}{{0x0}}%
\regfield{(reserved)}{1}{27}{0}%
\regfield{SDHOST\_RX\_WMARK}{11}{16}{{0x0}}%
\regfield{(reserved)}{4}{12}{0000}%
\regfield{SDHOST\_TX\_WMARK}{12}{0}{{0x0}}%
\reglabel{Reset}\regnewline%
\begin{regdesc}\begin{reglist}
\label{fielddesc:SDHOSTTXWMARK}\item [SDHOST\_TX\_WMARK] 配置向卡传输数据时的 FIFO 阈值。当 FIFO 数据计数小于或等于此数字时，会产生 DMA/FIFO 请求。在数据包结束时，无论阈值编程如何，都会生成请求或中断。\\
在非 DMA 模式下，当使能发送 FIFO 阈值 (TXDR) 中断时，会生成中断而不是 DMA 请求。在数据包结束时，最后一次中断时，主机负责只填充所需的剩余字节（不在 FIFO 满或 CIU 完成数据传输后，因为 FIFO 可能不为空）。\\
在 DMA 模式下，如果最后一次传输小于突发大小，DMA 控制器会进行单周期传输，直到传输所需的字节。 (R/W)
\label{fielddesc:SDHOSTRXWMARK}\item [SDHOST\_RX\_WMARK] 配置从卡接收数据时的 FIFO 阈值。当 FIFO 数据计数达到大于此数字时，会产生 DMA/FIFO 请求。在数据包结束时，无论阈值编程如何，都会生成请求，以完成任何剩余的数据。\\
在非 DMA 模式下，当使能接收器 FIFO 阈值 (RXDR) 中断时，会生成中断而不是 DMA 请求。在数据包结束时，如果阈值编程大于任何剩余的数据，中断不会生成。主机有责任在看到数据传输完成中断时读取剩余的字节。\\
在 DMA 模式下，即使剩余的字节少于阈值，DMA 请求也会进行单次传输，以在设置数据传输完成中断之前清除任何剩余的字节。 (R/W)
\label{fielddesc:SDHOSTDMAMULTIPLETRANSACTIONSIZE}\item [SDHOST\_DMA\_MULTIPLE\_TRANSACTION\_SIZE] 配置多次突发传输的大小，应与 DMA 控制器的多次传输大小 SDHOST\_SRC/DEST\_MSIZE 编程相同。\\
0：1 字节传输\\
1：4 字节传输\\
2：8 字节传输\\
3：16 字节传输\\
4：32 字节传输\\
5：64 字节传输\\
6：128 字节传输\\
7：256 字节传输\\ (R/W)
\end{reglist}\end{regdesc}
\end{register}


\begin{register}{H}{SDHOST\_DEBNCE\_REG}{0x0064}\label{regdesc:SDHOSTDEBNCEREG}
\regfield{(reserved)}{8}{24}{00000000}%
\regfield{SDHOST\_DEBOUNCE\_COUNT}{24}{0}{{0x0000}}%
\reglabel{Reset}\regnewline%
\begin{regdesc}\begin{reglist}
\label{fielddesc:SDHOSTDEBOUNCECOUNT}\item [SDHOST\_DEBOUNCE\_COUNT] 配置防抖滤波器逻辑使用的主机时钟数量。典型的防抖时间是 5 \verb+~+ 25 毫秒，防止卡插入或移除时不稳定。 (R/W)
\end{reglist}\end{regdesc}
\end{register}


\begin{register}{H}{SDHOST\_USRID\_REG}{0x0068}\label{regdesc:SDHOSTUSRIDREG}
\regfield{SDHOST\_USRID\_REG}{32}{0}{{0x000000}}%
\reglabel{Reset}\regnewline%
\begin{regdesc}\begin{reglist}
\label{fielddesc:SDHOSTUSRIDREG}\item [SDHOST\_USRID\_REG] 配置用户身份标识，值由用户配置。用户也可将该寄存器用作暂存寄存器。 (R/W)
\end{reglist}\end{regdesc}
\end{register}


\begin{register}{H}{SDHOST\_VERID\_REG}{0x{}006C}\label{regdesc:VERIDREG}
\regfield{SDHOST\_VERSIONID\_REG}{32}{0}{{0x{}5432270A}}%
\reglabel{Reset}\regnewline%
\begin{regdesc}\begin{reglist}
\item [SDHOST\_VERSIONID\_REG] 表示硬件版本。固件也可以读取该寄存器。 (RO)
\end{reglist}\end{regdesc}
\end{register}


\begin{register}{H}{SDHOST\_HCON\_REG}{0x{}0070}\label{regdesc:HCONREG}
\regfield{(reserved)}{5}{27}{{0x{}0}}%
\regfield{(reserved)}{1}{26}{{0x{}0}}%
\regfield{SDHOST\_NUM\_CLK\_DIV\_REG}{2}{24}{{0x{}3}}%
\regfield{(resvrved)}{1}{23}{{0x{}1}}%
\regfield{SDHOST\_HOLD\_REG}{1}{22}{{0x{}1}}%
\regfield{SDHOST\_RAM\_INDISE\_REG}{1}{21}{{0x{}0}}%
\regfield{SDHOST\_DMA\_WIDTH\_REG}{3}{18}{{0x{}1}}%
\regfield{(reserved)}{2}{16}{{0x{}0}}%
\regfield{SDHOST\_ADDR\_WIDTH\_REG}{6}{10}{{0x{}13}}%
\regfield{SDHOST\_DATA\_WIDTH\_REG}{3}{7}{{0x{}1}}%
\regfield{SDHOST\_BUS\_TYPE\_REG}{1}{6}{{0x{}1}}%
\regfield{SDHOST\_CARD\_NUM\_REG}{5}{1}{{0x{}1}}%
\regfield{SDHOST\_CARD\_TYPE\_REG}{1}{0}{{0x{}1}}%
\reglabel{Reset}\regnewline%
\begin{regdesc}\begin{reglist}
\item [SDHOST\_NUM\_CLK\_DIV\_REG] 表示设计中有 4 个时钟分频器。(RO)
\item [SDHOST\_HOLD\_REG] 表示数据路径中有一个保持寄存器。(RO)
\item [SDHOST\_RAM\_INDISE\_REG] 表示 SDMMC 模块内部包含 RAM。(RO)
\item [SDHOST\_DMA\_WIDTH\_REG] 表示 DMA 数据宽度为 32。(RO)
\item [SDHOST\_ADDR\_WIDTH\_REG] 表示寄存器地址宽度为 32。(RO)
\item [SDHOST\_DATA\_WIDTH\_REG] 表示寄存器数据宽度为 32。(RO)
\item [SDHOST\_BUS\_TYPE\_REG] 表示寄存器配置为 APB 总线。(RO)
\item [SDHOST\_CARD\_NUM\_REG] 表示支持的卡数量为 2。(RO)
\item [SDHOST\_CARD\_TYPE\_REG] 表示硬件支持 SDIO 和 MMC。(RO)
\end{reglist}\end{regdesc}
\end{register}


\begin{register}{H}{SDHOST\_UHS\_REG}{0x0074}\label{regdesc:SDHOSTUHSREG}
\regfield{(reserved)}{14}{18}{00000000000000}%
\regfield{SDHOST\_DDR\_REG}{2}{16}{{0x0}}%
\regfield{(reserved)}{14}{2}{00000000000000}%
\regfield{SDHOST\_VOLT\_REG}{2}{0}{{0x0}}%
\reglabel{Reset}\regnewline%
\begin{regdesc}\begin{reglist}
\label{fielddesc:SDHOSTVOLTREG}\item [SDHOST\_VOLT\_REG] 配置卡是否使用高电压模式。每张卡一位，位 0 对应卡 0。\\
0：电压 3.3 V\\
1：电压 1.8 V\\ (R/W)
\label{fielddesc:SDHOSTDDRREG}\item [SDHOST\_DDR\_REG] 配置卡是否使用 DDR（双倍数据速率）模式。每张卡一位，位 0 对应卡 0。\\
0：非 DDR 模式\\
1：DDR 模式\\ (R/W)
\end{reglist}\end{regdesc}
\end{register}


\begin{register}{H}{SDHOST\_RST\_N\_REG}{0x0078}\label{regdesc:SDHOSTRSTNREG}
\regfield{(reserved)}{30}{2}{000000000000000000000000000000}%
\regfield{SDHOST\_CARD\_RESET}{2}{0}{{0x1}}%
\reglabel{Reset}\regnewline%
\begin{regdesc}\begin{reglist}
\label{fielddesc:SDHOSTCARDRESET}\item [SDHOST\_CARD\_RESET] 配置卡的硬件复位。每张卡一位，位 0 对应卡 0。卡复位后进入预空闲状态，需要重新初始化。\\
0：复位\\
1：活跃模式\\ (R/W)
\end{reglist}\end{regdesc}
\end{register}


\begin{register}{H}{SDHOST\_BMOD\_REG}{0x0080}\label{regdesc:SDHOSTBMODREG}
\regfield{(reserved)}{21}{11}{000000000000000000000}%
\regfield{SDHOST\_BMOD\_PBL}{3}{8}{{0x0}}%
\regfield{SDHOST\_BMOD\_DE}{1}{7}{{0}}%
\regfield{(reserved)}{5}{2}{00000}%
\regfield{SDHOST\_BMOD\_FB}{1}{1}{{0}}%
\regfield{SDHOST\_BMOD\_SWR}{1}{0}{{0}}%
\reglabel{Reset}\regnewline%
\begin{regdesc}\begin{reglist}
\label{fielddesc:SDHOSTBMODSWR}\item [SDHOST\_BMOD\_SWR] 配置是否复位所有 DMA 内部寄存器。该位在一个时钟周期后自动清除。\\
0：无效\\
1：重置\\ (R/W)
\label{fielddesc:SDHOSTBMODFB}\item [SDHOST\_BMOD\_FB] 配置 AHB 主接口是否用于固定长度突发传输。\\
0：使用 单次 和 INCR（未指定长度）突发传输操作\\
1：仅使用单次、INCR4、INCR8 或 INCR16 突发传输\\ (R/W)
\label{fielddesc:SDHOSTBMODDE}\item [SDHOST\_BMOD\_DE] 配置是否使能 DMA。\\
0：不使能\\
1：使能\\ (R/W)
\label{fielddesc:SDHOSTBMODPBL}\item [SDHOST\_BMOD\_PBL] 配置一次 DMA 传输的最大传输长度。\\
0：1 拍传输\\
1：4 拍传输\\
2：8 拍传输\\
3：16 拍传输\\
4：32 拍传输\\
5：64 拍传输\\
6：128 拍传输\\
7：256 拍传输\\
 (R/W)
\end{reglist}\end{regdesc}
\end{register}


\begin{register}{H}{SDHOST\_PLDMND\_REG}{0x0084}\label{regdesc:SDHOSTPLDMNDREG}
\regfield{SDHOST\_PLDMND\_PD}{32}{0}{{0x000000}}%
\reglabel{Reset}\regnewline%
\begin{regdesc}\begin{reglist}
\label{fielddesc:SDHOSTPLDMNDPD}\item [SDHOST\_PLDMND\_PD] 配置 DMA FSM 是否恢复获取描述符。如果描述符的 OWNER 位未置 1，FSM 将进入挂起状态，软件需要向此字段写入任何值以使 FSM 恢复获取描述符。 (WO)
\end{reglist}\end{regdesc}
\end{register}


\begin{register}{H}{SDHOST\_DBADDR\_REG}{0x0088}\label{regdesc:SDHOSTDBADDRREG}
\regfield{SDHOST\_DBADDR\_REG}{32}{0}{{0x000000}}%
\reglabel{Reset}\regnewline%
\begin{regdesc}\begin{reglist}
\label{fielddesc:SDHOSTDBADDRREG}\item [SDHOST\_DBADDR\_REG] 配置链表第一个描述符的基地址。位 [1:0] 被内部 DMA 内部忽略并视为全零。  (R/W)
\end{reglist}\end{regdesc}
\end{register}


\begin{register}{H}{SDHOST\_CARDTHRCTL\_REG}{0x0100}\label{regdesc:SDHOSTCARDTHRCTLREG}
\regfield{SDHOST\_CARDTHRESHOLD\_REG}{16}{16}{{0x00}}%
\regfield{(reserved)}{14}{2}{00000000000000}%
\regfield{SDHOST\_CARDCLRINTEN\_REG}{1}{1}{{0}}%
\regfield{SDHOST\_CARDRDTHREN\_REG}{1}{0}{{0}}%
\reglabel{Reset}\regnewline%
\begin{regdesc}\begin{reglist}
\label{fielddesc:SDHOSTCARDRDTHRENREG}\item [SDHOST\_CARDRDTHREN\_REG] 配置是否使能卡读取阈值。\\
0：不使能\\
1：使能\\ (R/W)
\label{fielddesc:SDHOSTCARDCLRINTENREG}\item [SDHOST\_CARDCLRINTEN\_REG] 配置是否使能忙碌清除中断生成。\\
0：不使能\\
1：使能\\ (R/W)
\label{fielddesc:SDHOSTCARDTHRESHOLDREG}\item [SDHOST\_CARDTHRESHOLD\_REG] 配置卡读取阈值。单位：字节。仅在 SDHOST\_CARDRDTHREN\_REG 为 1 时有效，需小于内部 FIFO 大小 512。  (R/W)
\end{reglist}\end{regdesc}
\end{register}


\begin{register}{H}{SDHOST\_EMMCDDR\_REG}{0x010C}\label{regdesc:SDHOSTEMMCDDRREG}
\regfield{(reserved)}{30}{2}{000000000000000000000000000000}%
\regfield{SDHOST\_HALFSTARTBIT\_REG}{2}{0}{{0x0}}%
\reglabel{Reset}\regnewline%
\begin{regdesc}\begin{reglist}
\label{fielddesc:SDHOSTHALFSTARTBITREG}\item [SDHOST\_HALFSTARTBIT\_REG] 配置开始位检测机制的开始位检测长度。每张卡一位，位 0 对应卡 0。eMMC4.5 及以上版本此位需为 1，SD 应用此位需为 0。对于 eMMC4.5，起始位检测模式可配置为：\\
0：完整周期\\
1：少于一个完整周期\\(R/W)
\end{reglist}\end{regdesc}
\end{register}


\begin{register}{H}{SDHOST\_ENSHIFT\_REG}{0x0110}\label{regdesc:SDHOSTENSHIFTREG}
\regfield{(reserved)}{28}{4}{0000000000000000000000000000}%
\regfield{SDHOST\_ENABLE\_SHIFT\_REG}{4}{0}{{0x0}}%
\reglabel{Reset}\regnewline%
\begin{regdesc}\begin{reglist}
\label{fielddesc:SDHOSTENABLESHIFTREG}\item [SDHOST\_ENABLE\_SHIFT\_REG] 配置设计中默认的相位偏移量。每张卡 2 位，位 1:0 对应卡 0。\\
0x0：默认相位偏移\\
0x1：偏移到下个上升沿\\
0x2：偏移到下个下降沿\\
0x3：保留\\ (R/W)
\end{reglist}\end{regdesc}
\end{register}


\begin{register}{H}{SDHOST\_CLK\_EDGE\_SEL\_REG}{0x0800}\label{regdesc:SDHOSTCLKEDGESELREG}
\regfield{(reserved)}{9}{23}{000000001}%
\regfield{SDHOST\_ESD\_MODE}{1}{22}{{0}}%
\regfield{SDHOST\_ESDIO\_MODE}{1}{21}{{0}}%
\regfield{(reserved)}{21}{0}{000100000001000000000}%
\reglabel{Reset}\regnewline%
\begin{regdesc}\begin{reglist}
\label{fielddesc:SDHOSTESDIOMODE}\item [SDHOST\_ESDIO\_MODE] 配置是否使能 eSDIO 模式。\\
0：不使能\\
1：使能\\ (R/W)
\label{fielddesc:SDHOSTESDMODE}\item [SDHOST\_ESD\_MODE] 配置是否使能 eSD 模式。\\
0：不使能\\
1：使能\\ (R/W)
\end{reglist}\end{regdesc}
\end{register}


\begin{register}{H}{SDHOST\_DLL\_CLK\_CONF\_REG}{0x0808}\label{regdesc:SDHOSTDLLCLKCONFREG}
\regfield{(reserved)}{11}{21}{00000000000}%
\regfield{SDHOST\_DLL\_CCLK\_IN\_SAM\_PHASE}{6}{15}{{0}}%
\regfield{SDHOST\_DLL\_CCLK\_IN\_DRV\_PHASE}{6}{9}{{0}}%
\regfield{SDHOST\_DLL\_CCLK\_IN\_SLF\_PHASE}{6}{3}{{0}}%
\regfield{(reserved)}{3}{0}{000}%
\reglabel{Reset}\regnewline%
\begin{regdesc}\begin{reglist}
\label{fielddesc:SDHOSTDLLCCLKINSLFPHASE}\item [SDHOST\_DLL\_CCLK\_IN\_SLF\_PHASE] 配置选择高频时钟时内部信号的时钟相位。单位是时钟周期的 1/64。  (R/W)
\label{fielddesc:SDHOSTDLLCCLKINDRVPHASE}\item [SDHOST\_DLL\_CCLK\_IN\_DRV\_PHASE] 配置选择高频时钟时输出信号的时钟相位。单位是时钟周期的 1/64。  (R/W)
\label{fielddesc:SDHOSTDLLCCLKINSAMPHASE}\item [SDHOST\_DLL\_CCLK\_IN\_SAM\_PHASE] 配置选择高频时钟时输入信号的时钟相位。单位是时钟周期的 1/64。  (R/W)
\end{reglist}\end{regdesc}
\end{register}


\begin{register}{H}{SDHOST\_INTMASK\_REG}{0x0024}\label{regdesc:SDHOSTINTMASKREG}
\regfield{(reserved)}{14}{18}{00000000000000}%
\regfield{SDHOST\_SDIO\_INT\_MASK}{2}{16}{{0x0}}%
\regfield{SDHOST\_INT\_MASK}{16}{0}{{0x00}}%
\reglabel{Reset}\regnewline%
\begin{regdesc}\begin{reglist}
\label{fielddesc:SDHOSTINTMASK}\item [SDHOST\_INT\_MASK] 写 1 到对应位使能对应中断。\\
位 15 (EBE)：结束位错误/无 CRC 错误中断\\
位 14 (ACD)：自动命令结束中断\\
位 13 (SBE/BCI)：接收起始位错误中断\\
位 12 (HLE)：硬件锁定写入错误中断\\
位 11 (FRUN)：FIFO 空/满错误中断\\
位 10 (HTO)：主机填充数据超时中断\\
位 9 (DRTO)：数据读取超时中断\\
位 8 (RTO)：卡响应超时中断\\
位 7 (DCRC)：数据 CRC 错误中断\\
位 6 (RCRC)：卡响应 CRC 错误中断\\
位 5 (RXDR)：接收 FIFO 数据请求中断\\
位 4 (TXDR)：发送 FIFO 数据请求中断\\
位 3 (DTO)：数据传输结束中断\\
位 2 (CMDD)：命令执行完毕中断\\
位 1 (RE)：卡响应错误中断\\
位 0 (CD)：检测到卡中断\\  (R/W)
\label{fielddesc:SDHOSTSDIOINTMASK}\item [SDHOST\_SDIO\_INT\_MASK] 写 1 到对应位使能来自对应卡的中断。每张卡一位，位 0 对应卡 0。  (R/W)
\end{reglist}\end{regdesc}
\end{register}


\begin{register}{H}{SDHOST\_MINTSTS\_REG}{0x0040}\label{regdesc:SDHOSTMINTSTSREG}
\regfield{(reserved)}{14}{18}{00000000000000}%
\regfield{SDHOST\_SDIO\_INTERRUPT\_MSK}{2}{16}{{0x0}}%
\regfield{SDHOST\_INT\_STATUS\_MSK}{16}{0}{{0x00}}%
\reglabel{Reset}\regnewline%
\begin{regdesc}\begin{reglist}
\label{fielddesc:SDHOSTINTSTATUSMSK}\item [SDHOST\_INT\_STATUS\_MSK] 中断的屏蔽状态。每一位对应一个中断：\\
位 15 (EBE)：结束位错误/无 CRC 错误中断\\
位 14 (ACD)：自动命令结束中断\\
位 13 (SBE/BCI)：接收起始位错误中断\\
位 12 (HLE)：硬件锁定写入错误中断\\
位 11 (FRUN)：FIFO 空/满错误中断\\
位 10 (HTO)：主机填充数据超时中断\\
位 9 (DRTO)：数据读取超时中断\\
位 8 (RTO)：卡响应超时中断\\
位 7 (DCRC)：数据 CRC 错误中断\\
位 6 (RCRC)：卡响应 CRC 错误中断\\
位 5 (RXDR)：接收 FIFO 数据请求中断\\
位 4 (TXDR)：发送 FIFO 数据请求中断\\
位 3 (DTO)：数据传输结束中断\\
位 2 (CMDD)：命令执行完毕中断\\
位 1 (RE)：卡响应错误中断\\
位 0 (CD)：检测到卡中断\\ (RO)
\label{fielddesc:SDHOSTSDIOINTERRUPTMSK}\item [SDHOST\_SDIO\_INTERRUPT\_MSK] 来自 SDIO 卡的中断屏蔽状态。每张卡一位，位 0 对应卡 0。 (RO)
\end{reglist}\end{regdesc}
\end{register}


\begin{register}{H}{SDHOST\_RINTSTS\_REG}{0x0044}\label{regdesc:SDHOSTRINTSTSREG}
\regfield{(reserved)}{14}{18}{00000000000000}%
\regfield{SDHOST\_SDIO\_INTERRUPT\_RAW}{2}{16}{{0x0}}%
\regfield{SDHOST\_INT\_STATUS\_RAW}{16}{0}{{0x00}}%
\reglabel{Reset}\regnewline%
\begin{regdesc}\begin{reglist}
\label{fielddesc:SDHOSTINTSTATUSRAW}\item [SDHOST\_INT\_STATUS\_RAW] 中断的原始状态。每一位对应一个中断：\\
位 15 (EBE)：结束位错误/无 CRC 错误中断\\
位 14 (ACD)：自动命令结束中断\\
位 13 (SBE/BCI)：接收起始位错误中断\\
位 12 (HLE)：硬件锁定写入错误中断\\
位 11 (FRUN)：FIFO 空/满错误中断\\
位 10 (HTO)：主机填充数据超时中断\\
位 9 (DRTO)：数据读取超时中断\\
位 8 (RTO)：卡响应超时中断\\
位 7 (DCRC)：数据 CRC 错误中断\\
位 6 (RCRC)：卡响应 CRC 错误中断\\
位 5 (RXDR)：接收 FIFO 数据请求中断\\
位 4 (TXDR)：发送 FIFO 数据请求中断\\
位 3 (DTO)：数据传输结束中断\\
位 2 (CMDD)：命令执行完毕中断\\
位 1 (RE)：卡响应错误中断\\
位 0 (CD)：检测到卡中断\\ (R/W1C)
\label{fielddesc:SDHOSTSDIOINTERRUPTRAW}\item [SDHOST\_SDIO\_INTERRUPT\_RAW] 来自 SDIO 卡的中断原始状态。每张卡一位，位 0 对应卡 0。 (R/W1C)
\end{reglist}\end{regdesc}
\end{register}

\begin{register}{H}{SDHOST\_IDSTS\_REG}{0x008C}\label{regdesc:SDHOSTIDSTSREG}
\regfield{(reserved)}{15}{17}{000000000000000}%
\regfield{SDHOST\_IDSTS\_FSM}{4}{13}{{0x0}}%
\regfield{SDHOST\_IDSTS\_FBE\_CODE}{3}{10}{{0x0}}%
\regfield{SDHOST\_IDSTS\_AIS}{1}{9}{{0}}%
\regfield{SDHOST\_IDSTS\_NIS}{1}{8}{{0}}%
\regfield{(reserved)}{2}{6}{00}%
\regfield{SDHOST\_IDSTS\_CES}{1}{5}{{0}}%
\regfield{SDHOST\_IDSTS\_DU}{1}{4}{{0}}%
\regfield{(reserved)}{1}{3}{0}%
\regfield{SDHOST\_IDSTS\_FBE}{1}{2}{{0}}%
\regfield{SDHOST\_IDSTS\_RI}{1}{1}{{0}}%
\regfield{SDHOST\_IDSTS\_TI}{1}{0}{{0}}%
\reglabel{Reset}\regnewline%
\begin{regdesc}\begin{reglist}
\label{fielddesc:SDHOSTIDSTSTI}\item [SDHOST\_IDSTS\_TI] IDSTS\_TI 中断的原始中断状态。 (R/W1C)
\label{fielddesc:SDHOSTIDSTSRI}\item [SDHOST\_IDSTS\_RI] IDSTS\_RI 中断的原始中断状态。 (R/W1C)
\label{fielddesc:SDHOSTIDSTSFBE}\item [SDHOST\_IDSTS\_FBE] IDSTS\_FBE 中断的原始中断状态。 (R/W1C)
\label{fielddesc:SDHOSTIDSTSDU}\item [SDHOST\_IDSTS\_DU] IDSTS\_DU 中断的原始中断状态。 (R/W1C)
\label{fielddesc:SDHOSTIDSTSCES}\item [SDHOST\_IDSTS\_CES] 卡错误汇总中断的原始中断状态，为 SDHOST\_RINTSTS\_REG 中以下中断原始状态的逻辑或运算结果：\\
EBE：结束位错误/无 CRC 错误中断\\
SBE/BCI：接收起始位错误中断\\
DRTO：数据读取超时中断\\
RTO：卡响应超时中断\\
DCRC：数据 CRC 错误中断\\
RCRC：卡响应 CRC 错误中断\\
RE：卡响应错误中断\\
此该位为 1 时，DMA 会在响应错误时中止。 (R/W1C)
\label{fielddesc:SDHOSTIDSTSNIS}\item [SDHOST\_IDSTS\_NIS] 正常中断汇总的原始中断状态，为 SDHOST\_IDSTS\_REG[1:0] 的逻辑或运算结果。这是一个粘性位，软件每次清除导致此字段被设置的相应位时，必须也清除此位。 (R/W1C)
\label{fielddesc:SDHOSTIDSTSAIS}\item [SDHOST\_IDSTS\_AIS] 异常中断汇总的原始中断状态，为 SDHOST\_IDSTS\_REG[2] 和 SDHOST\_IDSTS\_REG[4] 的逻辑或运算结果。这是一个粘性位，软件每次清除导致此字段被设置的相应位时，必须也清除此位。 (R/W1C)
\item [见下页...]
\end{reglist}\end{regdesc}
\end{register}

\addtocounter{Regfloat}{-1}
\begin{register}{H}{SDHOST\_IDSTS\_REG}{0x008C}
\begin{regdesc}\begin{reglist}
\item [接上页...]
\label{fielddesc:SDHOSTIDSTSFBECODE}\item [SDHOST\_IDSTS\_FBE\_CODE] 表示导致总线错误的错误类型，仅在 SDHOST\_IDSTS\_REG[2] 为 1 时有效。\\
1：在传输过程中接收到主机中止\\
2：在接收过程中接收到主机中止\\
其他值：保留\\ (RO)
\label{fielddesc:SDHOSTIDSTSFSM}\item [SDHOST\_IDSTS\_FSM] 表示 DMA FSM 当前状态。\\
0：DMA\_IDLE（空闲状态）\\
1：DMA\_SUSPEND（暂停状态）\\
2：DESC\_RD（描述符读取状态）\\
3：DESC\_CHK（描述符检查状态）\\
4：DMA\_RD\_REQ\_WAIT（等待读取数据请求状态）\\
5：DMA\_WR\_REQ\_WAIT（等待写入数据请求状态）\\
6：DMA\_RD（数据读取状态）\\
7：DMA\_WR（数据写入状态）\\
8：DESC\_CLOSE（描述符关闭状态）\\ (RO)
\end{reglist}\end{regdesc}
\end{register}


\begin{register}{H}{SDHOST\_IDINTEN\_REG}{0x0090}\label{regdesc:SDHOSTIDINTENREG}
\regfield{(reserved)}{22}{10}{0000000000000000000000}%
\regfield{SDHOST\_IDINTEN\_AI}{1}{9}{{0}}%
\regfield{SDHOST\_IDINTEN\_NI}{1}{8}{{0}}%
\regfield{(reserved)}{2}{6}{00}%
\regfield{SDHOST\_IDINTEN\_CES}{1}{5}{{0}}%
\regfield{SDHOST\_IDINTEN\_DU}{1}{4}{{0}}%
\regfield{(reserved)}{1}{3}{0}%
\regfield{SDHOST\_IDINTEN\_FBE}{1}{2}{{0}}%
\regfield{SDHOST\_IDINTEN\_RI}{1}{1}{{0}}%
\regfield{SDHOST\_IDINTEN\_TI}{1}{0}{{0}}%
\reglabel{Reset}\regnewline%
\begin{regdesc}\begin{reglist}
\label{fielddesc:SDHOSTIDINTENTI}\item [SDHOST\_IDINTEN\_TI] 写入 1 使能 IDSTS\_TI 中断。 (R/W)
\label{fielddesc:SDHOSTIDINTENRI}\item [SDHOST\_IDINTEN\_RI] 写入 1 使能 IDSTS\_RI 中断。  (R/W)
\label{fielddesc:SDHOSTIDINTENFBE}\item [SDHOST\_IDINTEN\_FBE] 写入 1 使能 IDSTS\_FBE 中断。  (R/W)
\label{fielddesc:SDHOSTIDINTENDU}\item [SDHOST\_IDINTEN\_DU] 写入 1 使能 IDSTS\_DU 中断。  (R/W)
\label{fielddesc:SDHOSTIDINTENCES}\item [SDHOST\_IDINTEN\_CES] 写入 1 使能卡错误汇总中断。  (R/W)
\label{fielddesc:SDHOSTIDINTENNI}\item [SDHOST\_IDINTEN\_NI] 写入 1 使能正常中断汇总。  (R/W)
\label{fielddesc:SDHOSTIDINTENAI}\item [SDHOST\_IDINTEN\_AI] 写入 1 使能异常中断汇总。  (R/W)
\end{reglist}\end{regdesc}
\end{register}


\begin{register}{H}{SDHOST\_RESP0\_REG}{0x0030}\label{regdesc:SDHOSTRESP0REG}
\regfield{SDHOST\_RESPONSE0\_REG}{32}{0}{{0x000000}}%
\reglabel{Reset}\regnewline%
\begin{regdesc}\begin{reglist}
\label{fielddesc:SDHOSTRESPONSE0REG}\item [SDHOST\_RESPONSE0\_REG] 表示响应的位 [31:0]。 (RO)
\end{reglist}\end{regdesc}
\end{register}


\begin{register}{H}{SDHOST\_RESP1\_REG}{0x0034}\label{regdesc:SDHOSTRESP1REG}
\regfield{SDHOST\_RESPONSE1\_REG}{32}{0}{{0x000000}}%
\reglabel{Reset}\regnewline%
\begin{regdesc}\begin{reglist}
\label{fielddesc:SDHOSTRESPONSE1REG}\item [SDHOST\_RESPONSE1\_REG] 表示长响应的位 [63:32]。 (RO)
\end{reglist}\end{regdesc}
\end{register}


\begin{register}{H}{SDHOST\_RESP2\_REG}{0x0038}\label{regdesc:SDHOSTRESP2REG}
\regfield{SDHOST\_RESPONSE2\_REG}{32}{0}{{0x000000}}%
\reglabel{Reset}\regnewline%
\begin{regdesc}\begin{reglist}
\label{fielddesc:SDHOSTRESPONSE2REG}\item [SDHOST\_RESPONSE2\_REG] 表示长响应的位 [95:64]。 (RO)
\end{reglist}\end{regdesc}
\end{register}


\begin{register}{H}{SDHOST\_RESP3\_REG}{0x003C}\label{regdesc:SDHOSTRESP3REG}
\regfield{SDHOST\_RESPONSE3\_REG}{32}{0}{{0x000000}}%
\reglabel{Reset}\regnewline%
\begin{regdesc}\begin{reglist}
\label{fielddesc:SDHOSTRESPONSE3REG}\item [SDHOST\_RESPONSE3\_REG] 表示长响应的位 [127:96]。 (RO)
\end{reglist}\end{regdesc}
\end{register}


\begin{register}{H}{SDHOST\_STATUS\_REG}{0x0048}\label{regdesc:SDHOSTSTATUSREG}
\regfield{(reserved)}{2}{30}{00}%
\regfield{SDHOST\_FIFO\_COUNT}{13}{17}{{0x00}}%
\regfield{SDHOST\_RESPONSE\_INDEX}{6}{11}{{0x0}}%
\regfield{SDHOST\_DATA\_STATE\_MC\_BUSY}{1}{10}{{1}}%
\regfield{SDHOST\_DATA\_BUSY}{1}{9}{{1}}%
\regfield{SDHOST\_DATA\_3\_STATUS}{1}{8}{{1}}%
\regfield{SDHOST\_COMMAND\_FSM\_STATES}{4}{4}{{0x1}}%
\regfield{SDHOST\_FIFO\_FULL}{1}{3}{{0}}%
\regfield{SDHOST\_FIFO\_EMPTY}{1}{2}{{1}}%
\regfield{SDHOST\_FIFO\_TX\_WATERMARK}{1}{1}{{1}}%
\regfield{SDHOST\_FIFO\_RX\_WATERMARK}{1}{0}{{0}}%
\reglabel{Reset}\regnewline%
\begin{regdesc}\begin{reglist}
\label{fielddesc:SDHOSTFIFORXWATERMARK}\item [SDHOST\_FIFO\_RX\_WATERMARK] 表示 FIFO 是否达到接收阈值，不关注数据是否传输。 \\
0：未达到\\
1：已达到\\(RO)
\label{fielddesc:SDHOSTFIFOTXWATERMARK}\item [SDHOST\_FIFO\_TX\_WATERMARK] 表示 FIFO 是否达到发送阈值，不关注数据是否传输。 \\
0：未达到\\
1：已达到\\(RO)
\label{fielddesc:SDHOSTFIFOEMPTY}\item [SDHOST\_FIFO\_EMPTY] 表示 FIFO 是否已空。\\
0：未空\\
1：已空\\(RO)
\label{fielddesc:SDHOSTFIFOFULL}\item [SDHOST\_FIFO\_FULL] 表示 FIFO 是否已满。\\
0：未满\\
1：已满\\(RO)
\item [见下页...]
\end{reglist}\end{regdesc}
\end{register}

\addtocounter{Regfloat}{-1}
\begin{register}{H}{SDHOST\_STATUS\_REG}{0x0048}
\begin{regdesc}\begin{reglist}
\item [接上页...]
\label{fielddesc:SDHOSTCOMMANDFSMSTATES}\item [SDHOST\_COMMAND\_FSM\_STATES] 表示命令 FSM 状态。\\
0x0：空闲\\
0x1：发送初始化序列\\
0x2：发送命令开始位\\
0x3：发送命令接收位\\
0x4：发送命令索引编号和参数\\
0x5：发送命令 CRC-7\\
0x6：发送命令结束位\\
0x7：接收响应开始位\\
0x8：接收响应 IRQ 响应\\
0x9：接收响应发送位\\
0xA：接收响应命令索引编号\\
0xB：接收响应数据\\
0xC：接收响应 CRC-7\\
0xD：接收响应结束位\\
0xE：命令通路等待 NCC\\
0xF：等待，命令到响应转换\\(RO)
\label{fielddesc:SDHOSTDATA3STATUS}\item [SDHOST\_DATA\_3\_STATUS] 表示选择的 sdhost\_card\_data[3] 输入信号值，检查是否存在卡。\\
0：卡不存在\\
1：卡存在\\ (RO)
\label{fielddesc:SDHOSTDATABUSY}\item [SDHOST\_DATA\_BUSY] 表示选择的 sdhost\_card\_data[0] 输入信号反向值，检查卡数据是否忙碌。\\
0：卡数据不忙碌\\
1：卡数据忙碌\\ (RO)
\label{fielddesc:SDHOSTDATASTATEMCBUSY}\item [SDHOST\_DATA\_STATE\_MC\_BUSY] 表示数据传输或接收状态机是否忙碌。\\
0：不忙碌\\
1：忙碌\\(RO)
\label{fielddesc:SDHOSTRESPONSEINDEX}\item [SDHOST\_RESPONSE\_INDEX] 表示前一个响应的索引，包括由核心发送的任何自动停止。 (RO)
\label{fielddesc:SDHOSTFIFOCOUNT}\item [SDHOST\_FIFO\_COUNT] 表示 FIFO 中已填充位置的数量。 (RO)
\end{reglist}\end{regdesc}
\end{register}


\begin{register}{H}{SDHOST\_CDETECT\_REG}{0x0050}\label{regdesc:SDHOSTCDETECTREG}
\regfield{(reserved)}{30}{2}{000000000000000000000000000000}%
\regfield{SDHOST\_CARD\_DETECT\_N}{2}{0}{{0x0}}%
\reglabel{Reset}\regnewline%
\begin{regdesc}\begin{reglist}
\label{fielddesc:SDHOSTCARDDETECTN}\item [SDHOST\_CARD\_DETECT\_N] 表示 sdhost\_card\_detect\_n 输入信号的值。每张卡一位，位 0 对应卡 0。 (RO)
\end{reglist}\end{regdesc}
\end{register}


\begin{register}{H}{SDHOST\_WRTPRT\_REG}{0x0054}\label{regdesc:SDHOSTWRTPRTREG}
\regfield{(reserved)}{30}{2}{000000000000000000000000000000}%
\regfield{SDHOST\_WRITE\_PROTECT}{2}{0}{{0x0}}%
\reglabel{Reset}\regnewline%
\begin{regdesc}\begin{reglist}
\label{fielddesc:SDHOSTWRITEPROTECT}\item [SDHOST\_WRITE\_PROTECT] 表示 sdhost\_card\_write\_prt 输入信号上的值，每张卡一位，位 0 对应卡 0。 (RO)
\end{reglist}\end{regdesc}
\end{register}


\begin{register}{H}{SDHOST\_TCBCNT\_REG}{0x005C}\label{regdesc:SDHOSTTCBCNTREG}
\regfield{SDHOST\_TCBCNT\_REG}{32}{0}{{0x000000}}%
\reglabel{Reset}\regnewline%
\begin{regdesc}\begin{reglist}
\label{fielddesc:SDHOSTTCBCNTREG}\item [SDHOST\_TCBCNT\_REG] 表示 CIU 传输到卡的字节数。 (RO)
\end{reglist}\end{regdesc}
\end{register}


\begin{register}{H}{SDHOST\_TBBCNT\_REG}{0x0060}\label{regdesc:SDHOSTTBBCNTREG}
\regfield{SDHOST\_TBBCNT\_REG}{32}{0}{{0x000000}}%
\reglabel{Reset}\regnewline%
\begin{regdesc}\begin{reglist}
\label{fielddesc:SDHOSTTBBCNTREG}\item [SDHOST\_TBBCNT\_REG] 表示主机或 DMA 内存与 BIU FIFO 之间传输的字节数。 (RO)
\end{reglist}\end{regdesc}
\end{register}


\begin{register}{H}{SDHOST\_DSCADDR\_REG}{0x0094}\label{regdesc:SDHOSTDSCADDRREG}
\regfield{SDHOST\_DSCADDR\_REG}{32}{0}{{0x000000}}%
\reglabel{Reset}\regnewline%
\begin{regdesc}\begin{reglist}
\label{fielddesc:SDHOSTDSCADDRREG}\item [SDHOST\_DSCADDR\_REG] 表示 DMA 当前读取的描述符起始地址。 (RO)
\end{reglist}\end{regdesc}
\end{register}


\begin{register}{H}{SDHOST\_BUFADDR\_REG}{0x0098}\label{regdesc:SDHOSTBUFADDRREG}
\regfield{SDHOST\_BUFADDR\_REG}{32}{0}{{0x000000}}%
\reglabel{Reset}\regnewline%
\begin{regdesc}\begin{reglist}
\label{fielddesc:SDHOSTBUFADDRREG}\item [SDHOST\_BUFADDR\_REG] 表示 DMA 当前访问的数据缓存地址。 (RO)
\end{reglist}\end{regdesc}
\end{register}


\begin{register}{H}{SDHOST\_BUFFIFO\_REG}{0x0200}\label{regdesc:SDHOSTBUFFIFOREG}
\regfield{SDHOST\_BUFFIFO\_REG}{32}{0}{{0x000000}}%
\reglabel{Reset}\regnewline%
\begin{regdesc}\begin{reglist}
\label{fielddesc:SDHOSTBUFFIFOREG}\item [SDHOST\_BUFFIFO\_REG] FIFO 的接受数据或待发送数据。  (R/W)
\end{reglist}\end{regdesc}
\end{register}
